%% ---------------------------------------------------------------------
%%  FIGURA - Maquina de estados da partida (DG-04)
%% ---------------------------------------------------------------------
\begin{figure}[htb]
	\caption{Máquina de estados da partida e transições}
	\label{fig:fsm}
	\centering
	\begin{tikzpicture}[
			font=\footnotesize,
			estado/.style={draw, rounded corners=3pt, minimum width=2.4cm, minimum height=0.85cm, align=center},
			inicial/.style={estado, very thick, fill=black!8},
			>=Stealth,
			tr/.style={->, thick},
			rot/.style={font=\scriptsize, inner sep=2pt, fill=white}
		]

		\node[inicial] (ready) at (0,0) {\texttt{ready}};
		\node[estado] (run) at (4.6,0) {\texttt{running}};
		\node[estado] (pause) at (4.6,-2.3) {\texttt{paused}};
		\node[estado] (over) at (9.2,1.1) {\texttt{over}};
		\node[estado] (won) at (9.2,-1.1) {\texttt{won}};

		\draw[tr] (ready) -- node[rot, above] {direção válida} (run);
		\draw[tr] (run) -- node[rot, right] {pausa ou perda de foco} (pause);
		\draw[tr] (pause) -- node[rot, left] {retomada} (run);
		\draw[tr] (run) -- node[rot, pos=0.62, above] {colisão} (over);
		\draw[tr] (run) -- node[rot, pos=0.62, below] {grade completa} (won);

		\draw[tr] (run) to[out=110, in=70, looseness=7] node[rot, above] {passo de simulação} (run);

		\draw[tr] (over) to[out=200, in=30, looseness=1.3] node[rot, pos=0.42, above] {reinício} (ready);
		\draw[tr] (won) to[out=200, in=-10, looseness=1.15] node[rot, pos=0.4, below] {reinício} (ready);
	\end{tikzpicture}
	\legend{Fonte: elaborado pelo autor (2026) a partir de \texttt{index.html}, seção~7.
	O estado \texttt{won} não constava da solicitação original: ele foi criado para tratar o caso de
	borda ``grade completamente ocupada''.}
\end{figure}
